\chapter{Stage one: foundations}

\section{Why it matters}

Every part of \method{} is a neural network trained by gradient descent. Without a working mental model of
how a loss produces gradients, how an optimiser uses them and why normalisation and residual connections make
deep networks trainable, the rest of the project is a set of incantations. This stage builds that model from
scratch.

\section{What to learn}

\begin{itemize}
  \item Linear classifiers, softmax and cross-entropy; what a loss surface is.
  \item Backpropagation as the chain rule on a computation graph, implemented by hand once.
  \item Stochastic gradient descent, momentum and Adam~\cite{kingma2015adam}; learning-rate schedules and warm-up.
  \item Convolutions, padding, stride and receptive fields; pooling; feature maps as images of features.
  \item Batch normalisation~\cite{ioffe2015batchnorm} and why it stabilises training.
  \item Residual connections~\cite{he2016resnet}: why a block that starts as the identity is easy to optimise.
        \method{} uses the same trick, starting every fusion block as an average.
  \item Overfitting, data augmentation, train, validation and test splits.
  \item Attention and the Transformer~\cite{vaswani2017attention}, and its application to images~\cite{dosovitskiy2021vit}.
  \item Mixed precision: what float16 can and cannot represent. The project hit this directly, as Chapter~\ref{ch:research} explains.
\end{itemize}

\section{Lectures and reading}

\begin{tabularx}{\textwidth}{@{}P{5.2cm}L@{}}
\hd{resource} & \hd{link and what to take from it} \\
Stanford CS231n, course site and notes & \link{https://cs231n.stanford.edu/} \newline \link{https://cs231n.github.io/} \newline The notes on optimisation, backpropagation and convolutional networks are the core text. \\
CS231n lectures, Spring 2017 & \link{https://www.youtube.com/playlist?list=PL3FW7Lu3i5JvHM8ljYj-zLfQRF3EO8sYv} \newline Lectures 1 to 11. Do assignments 1 and 2. \\
Neural Networks: Zero to Hero, Andrej Karpathy & \link{https://karpathy.ai/zero-to-hero.html} \newline Builds backpropagation and a training loop from nothing; watch alongside CS231n. \\
Deep Learning for Computer Vision, Michigan EECS 498 & \link{https://www.youtube.com/playlist?list=PL5-TkQAfAZFbzxjBHtzdVCWE0Zbhomg7r} \newline A more recent course by Justin Johnson, with an excellent detection lecture. \\
3Blue1Brown, neural networks & \link{https://www.3blue1brown.com/topics/neural-networks} \newline Geometric intuition for gradient descent and backpropagation. \\
Dive into Deep Learning & \link{https://d2l.ai/} \newline A free book with runnable code for every chapter; use it as a reference. \\
\end{tabularx}

\section{Papers}

\begin{tabularx}{\textwidth}{@{}P{5.2cm}L@{}}
\hd{paper} & \hd{take away} \\
ResNet~\cite{he2016resnet} & identity shortcuts make very deep networks trainable \\
Batch Normalization~\cite{ioffe2015batchnorm} & normalising activations stabilises and speeds up training \\
Adam~\cite{kingma2015adam} & per-parameter adaptive step sizes; the default optimiser here \\
Attention Is All You Need~\cite{vaswani2017attention} & attention as content-based mixing of a sequence \\
ViT~\cite{dosovitskiy2021vit} & images as sequences of patches; the bridge to sequence models for vision \\
\end{tabularx}

\section{In this repository}

Read \file{cffm-net-pilot/src/cffm/scan.py}. The reference scan is plain PyTorch and makes a good first real
code read: a recurrence, a chunked parallel version of it, and a test that the two agree.

\section{Exercises}

\begin{enumerate}
  \item Implement a two-layer network and its backward pass in NumPy; check every gradient numerically.
  \item Train a small convolutional network on CIFAR-10 and plot training and validation loss. Make it overfit
        on purpose, then fix it with augmentation.
  \item Remove the shortcut from a twenty-layer residual network and compare the training curves.
  \item In float16, compute $0.9^{1000}$ by repeated multiplication. Explain the result.
\end{enumerate}

\section{You are ready when}

You can explain, without notes, why a residual block that starts as the identity trains more easily than a
plain block, and you can write a correct training loop for an image classifier from an empty file.
